\documentclass[11pt]{article}
\usepackage[T1]{fontenc}
\usepackage{lmodern,microtype}
\usepackage[a4paper,margin=27mm]{geometry}
\usepackage{amsmath,amssymb,amsthm,mathtools,booktabs,longtable,array}
\usepackage[hidelinks]{hyperref}
\usepackage{enumitem}
\setlist{itemsep=2pt,topsep=4pt}
\allowdisplaybreaks
\setlength{\emergencystretch}{3em}
\newtheorem{theorem}{Theorem}[section]
\newtheorem{proposition}[theorem]{Proposition}
\newtheorem{lemma}[theorem]{Lemma}
\newtheorem{corollary}[theorem]{Corollary}
\theoremstyle{definition}
\newtheorem{definition}[theorem]{Definition}
\theoremstyle{remark}
\newtheorem{remark}[theorem]{Remark}
\newcommand{\CC}{\mathbb C}
\newcommand{\QQ}{\mathbb Q}
\newcommand{\ZZ}{\mathbb Z}
\newcommand{\FF}{\mathbb F}
\newcommand{\one}{\mathbf1}
\newcommand{\FPdim}{\operatorname{FPdim}}
\newcommand{\Irr}{\operatorname{Irr}}
\newcommand{\Rep}{\operatorname{Rep}}
\renewcommand{\Vec}{\operatorname{Vec}}
\newcommand{\Tr}{\operatorname{Tr}}
\newcommand{\supp}{\operatorname{supp}}
\newcommand{\rk}{\operatorname{rank}}
\newcommand{\diag}{\operatorname{diag}}
\newcommand{\RR}[2]{R^{(6)}_{#1,#2}}
\newcommand{\filename}[1]{\nolinkurl{#1}}
\numberwithin{equation}{section}
\newcommand{\NegativeCount}{16536}
\newcommand{\UnitaryCount}{53}
\newcommand{\SphericalOnlyCount}{1}
\newcommand{\PositiveCount}{54}
\newcommand{\UnresolvedCount}{22}
\newcommand{\CompleteBound}{3}
\newcommand{\CharacterExampleCount}{23}
\newcommand{\CensusCount}{16612}
\newcommand{\CensusBound}{8}
\newcommand{\ConstructorCount}{61}

\hypersetup{pdftitle={Verification guide: rank-six categorification through multiplicity nine},pdfauthor={}}
\title{Verification guide\\\large Rank-six categorification through multiplicity nine}
\author{}\date{17 October 2026}
\begin{document}\maketitle
\paragraph{Research status.} \textbf{AI-generated research \textperiodcentered{} independent verification pending.} Treat this manuscript and its accompanying research data with caution: they have not been independently verified by a human mathematician or formalized in Lean.

\section{Scope and authoritative records}
The package classifies categorifiability of the $\CensusCount$ rank-six
fusion rings of multiplicity at most nine in the published
\emph{Fusion Ring Census}. The current ledger has $\NegativeCount$
general exclusions, $\UnitaryCount$ unitary-positive rings, one
spherical-positive ring admitting no pseudo-unitary categorification,
and $\UnresolvedCount$ unresolved rings. Positive based rings, rather
than the number of constructions, are counted: the $\PositiveCount$
positive rings have $\ConstructorCount$ constructor records.

The authoritative files are \filename{data/census.json},
\filename{data/decisions.json} and
\filename{data/classification_counts.json}. The CSV ledger and paper
tables are derived from them. The model catalogue records the source,
categorical structure and complete basis permutation for each model.
The file \filename{data/unresolved_tensors.json} contains all
complete remaining tensors. General, pivotal, spherical,
pseudo-unitary and unitary existence are separate fields.

The census is pinned to
\href{https://github.com/sebastienpalcoux/fusion-ring-census}{\emph{Fusion Ring Census}},
revision
\begin{center}\small\filename{2a340871242879665fa7ba2b01fa15b8bdeaada3}.\end{center}
Its rank-six exact-multiplicity counts are
$39,154,384,872,1582,2768,3814,6999,8809$.
The compressed tensor file \filename{data/rank6_tables.txt.gz} has
SHA-256
\begin{center}\small
\filename{0e7552e17620452f237e94751c40cd93b0757a466168629db8cc1bf0b9451d00}.
\end{center}
The source and both compressed and uncompressed hashes are recorded in
\filename{provenance/repository_source.json}. Import checks fix the
input correspondence; completeness of the counting input is supplied
by the cited census. The classification does not rerun census enumeration.

\section{Verification scope and reproducible commands}
The completed verification campaign covers all $25\,362$ general-exclusion
certificates using the unchanged mathematical checker. Its full initial
invocation records $25\,355$ passes and seven failures, with exit status
one. Complete replays of all seven failed certificates then pass with
SymPy's Python ground types selected before import. The two post-loop
count and positive-registry checks also pass when run separately. The
initial invocation completed the census import, positive-model and
pre-negative checks. The aggregate receipt binds every accepted
certificate to its source and input hashes and checks complete coverage.

The seven selected certificates are $\RR{4}{416}$, $\RR{4}{427}$,
$\RR{4}{446}$, $\RR{5}{638}$, $\RR{6}{1035}$, $\RR{8}{311}$ and
$\RR{8}{3664}$. Their saved separating vectors use a positional ordering
of exact balance equations. Different string representations of
algebraic coefficients in SymPy's Python and Flint backends can change
that ordering and its row signs. The diagnostic packets verify the
corresponding exact signed permutations. Neither the mathematical
checker nor the certificate vectors were relaxed or replaced.

This is a combined campaign, not a single successful full invocation or
an all-Python full-census replay. The initial invocation did not record
its loaded backend and version at launch; the selected replays record
theirs explicitly. The complete scope and preserved failures are in
\filename{evidence/prepublication_20261017/rank6_backend/aggregate_summary.json}.
The accompanying coverage file contains one record for every excluded
ring. These checks do not provide independent human verification or a
formal proof in Lean.

After restoring the repository payloads and installing the dependencies
in \filename{supplements/rank6/requirements.txt}, the supplied orchestration
entry point is, from the repository root:
\begin{verbatim}
python evidence/prepublication_20261017/rank6_backend/replay_complete.py \
  --repo . --output-dir /tmp/rank6-fresh-campaign --jobs 2
\end{verbatim}
The output directory must not already exist. C++17 is required for the
finite-search helpers. The orchestrator records actual process
backends, preserves failures, repeats every failed certificate with
Python ground types, runs the post-loop checks and verifies complete
coverage. Its identity, coverage and rejection controls have been
tested; the orchestrator itself has not been executed end to end.
The completed campaign used the separately recorded full invocation
and selected replay commands.

Individual package checks, from \filename{supplements/rank6/}, include
\begin{verbatim}
python check_manifest.py
python software/current_regressions.py
python -O software/conformal_models.py --verify
python -O software/replay_m9.py --jobs 2
python -O software/replay_m9_structure.py
python -O software/replay_m9_completion.py
python software/release_consistency.py
\end{verbatim}
The manifest should be checked before generating new report files.
The optional \filename{--proofs-only} mode of \filename{verify.py}
omits positive reconstruction and is not a complete release check.
The structure-only replay concerns its specified categorical structures;
it cannot establish unrestricted nonexistence.

A negative checker reconstructs number fields, finite branch domains,
signed quotients, separating vectors and exact inequalities from the
full tensor. Numerical optimizers may propose certificates during
discovery; their entry points are disabled during certificate checking.
An incomplete search, timeout or numerical infeasibility flag cannot
establish an exclusion. A successful necessary system cannot establish
existence. Normal and optimized runs of the continuation checks verify
that their correctness checks do not depend on removable assertions.
The categorical theorems and the proofs of their application remain
mathematical inputs. Checksums establish file integrity, not a
mathematical conclusion.

\paragraph{The multiplicity-five completion.}
The exclusion of $\RR{5}{1204}$ first reconstructs its complete
$171$-profile domain. A forced central product removes $72$ profiles;
integral cyclic-rotation multiplicities remove the other $99$ using
$112$ certified integer bounds and $99$ final exact separating rays.
The proof and solver-disabled replay are under
\filename{research/exclusions/R6_5_1204/}:
\begin{verbatim}
python -O research/exclusions/R6_5_1204/verify_exclusion.py
\end{verbatim}
The inherited finite-profile certificate remains part of the proof.

\section{Completion at multiplicity nine}
Every one of the $8\,809$ rings of exact multiplicity nine is generally
excluded. The eleven final exclusions combine signed dimension-field and
different-ideal restrictions, diagonal associativity and trace saturation,
complete Frobenius--Schur indicator systems, and a Galois spin-orbit bound.
The proof does not assume the pivotality conjecture.

The certificate ledger is \filename{certificates/m9_completion.jsonl}.
Portable proof bundles are in
\filename{research/exclusions/m9_completion/}; each contains its complete
input tensors, mathematical argument and exact finite replay. Run
\begin{verbatim}
python software/replay_m9_completion.py
python -O software/replay_m9_completion.py
\end{verbatim}
The source-bound receipt is recorded in
\filename{verification/m9/completion/}. The checker regenerates the
required finite domains, verifies every branch and rejects corrupted
certificates. All correctness checks remain active with \filename{-O}.
This calculation concerns categorification; its counting input is the cited census.

\section{The three conformal constructions}
The positive registry \filename{data/conformal_positive_models.json}
adds $\RR{4}{22}$, $\RR{4}{91}$ and $\RR{4}{629}$.
Their models are the neutral components of the module categories for
$(B_2)_7\subset(D_7)_1$, $(C_3)_4\subset(B_{10})_1$ and
$(B_2)_{12}\subset(E_8)_1$, respectively.
The source files, exact certificates and independent mathematical audit
are organized under \filename{research/structural_models/}.

Each conformal embedding supplies an actual unitary algebra and module
category. The parity grading gives a rank-six neutral component.
Published action graphs or exact free-induction Gram matrices determine
a small set of possible fusion tensors. Local modules, duality,
formal-codegree induction and the known center's dimensions or
indicators distinguish the required tensor. Every resulting tensor is
compared in all $216$ entries with the census representative.
The proofs do not infer a category from dimensions, a modular invariant,
or matching modular data.

The canonical command
\begin{verbatim}
python -O software/conformal_models.py --verify
\end{verbatim}
checks the source bindings and exact tensor identifications. The
per-model directories \filename{quartic/}, \filename{c3_model/} and
\filename{r629/} contain the reconstruction programs and mathematical
proofs, including the deeper graph, Weyl-character and indicator checks.
These constructions prove existence. They do not classify all
associators on the three rings or identify every archived modular-data
matrix in a prescribed labelling.

\section{The remaining problems}
The canonical index \filename{research/frontier/} describes the two
remaining general existence questions, their necessary reductions and
exact replay commands. No multiplicity-nine case remains open.

The two residuals through multiplicity eight satisfy the
following necessary conditions. For $\RR{4}{155}$, the generalized near-group fusion rules give a
uniform sphericality and Galois-pseudo-unitarity proof by one norm
argument.
Fermion fixed-point traces and odd Gauss sums leave one necessary
induction profile, of center rank $52$, whose twist spectrum is forced
and has exponent $24$. The short exact replay is
\filename{research/frontier/R6_4_155/verify_center_cuts.py}.
Its pointed center also rules out every braided representative of the
Morita class. The exact reconstruction determines every entry of the
rank-$26$ centralizer's modular datum. Its Verlinde coefficients are
nonnegative integers, with maximum $9$. It forces ten self-dual
objects in that factor and twenty in the full center. All stated finite
ordinary and generalized indicator checks and all ordered rank-$26$
Vafa equations pass. No categorical construction follows.

For $\RR{8}{1660}$, the Ising subcategory makes its complement a
regular bimodule. Its relative center has rank eighteen and rules
$B[t]/(t^2-At-L)$, where $B$ is the Ising-center ring.
Formal-codegree divisibility and dimension-field constraints give
$58$ labelled choices. A central second-indicator obstruction and a
Morita-invariant invertible count leave one based ring, up to exchanging
the Ising factors.
The proof and exact replay are in
\filename{research/frontier/R6_8_1660/relative_center/}.
The retained full-center certificates force exponent dividing $48$.
Fermion orbit constraints and a weighted second-indicator congruence
reduce their $1764$ labelled profiles to $1727$. The tensor involution and
homogeneous-packet Vafa equations then leave $1670$; the exact replays are in
\filename{research/frontier/R6_8_1660/fermion_profiles/} and
\filename{research/frontier/R6_8_1660/spin_profiles/}.
A subsequent signed Galois--fermion four-divisibility condition removes
raw profile $133$ in low frame zero, leaving $1669$ profiles in branches
$163,753,753$. Independently, every indecomposable module category has
six simples; full-center saturation, integral Morita Gram reconstruction
and a weighted second-indicator obstruction leave three possible dual
based rings, of ranks six, six and eight.

The exact calculations and proof audits are under
\filename{research/frontier/new_research/}. From the package root, run
\begin{verbatim}
python research/frontier/new_research/replay.py
\end{verbatim}
The summaries of this finite system are in
\filename{data/frontier_research_progress.json}.

\subsection{Further necessary conditions}
The manuscript includes the graded-packet identities,
tensor-quotient obstruction, restrictions on Tannakian subcategories and
braided Morita models, and exact scope of the analytic tests.
The graded-packet conditions leave $1136=120+508+508$ profiles.
All-order indicator identities determine the exact central phase masses;
the full phase Gram constraints leave $640=68+286+286$ profiles.
Exactly eight spin simples have restriction $2X+2gX+\sigma+3T$.
For $\RR{4}{155}$, odd cyclic Clifford factorization is a necessary
condition at every odd order; all $1144$ tested central-target rotation
blocks at orders $3,5,\ldots,23$ pass. The all-order path-dimension
condition is equivalent to even multiplicity in the stated ring family.
Both $\RR{4}{155}$ and $\RR{8}{1660}$ remain unresolved.

The source-bound receipts and verifiers for these results are retained
at the repository path
\filename{evidence/residual_20261002/rank6/};
the current research summary and all-order refinements are recorded at
\filename{evidence/rigidity_20261017/rank6/}, with its
\filename{research_progress.json}.
The exact all-exponent Minkowski couplings are retained under
\filename{evidence/rank6_laptop_20261002/all_p_minkowski/}.
The source hashes in these records identify the inputs required by each
checker.

The manuscript further gives a theoretical proof that every tensor
functor from a hypothetical categorification of either residual ring
is fully faithful. It proves pseudo-unitarity of dominant images and
preservation of canonical pivotal structures, then combines these
facts with the quotient obstruction and the complete twist-one
restriction table. This strengthens the tensor-quotient restriction
without changing either unresolved verdict.

\subsection{All-order necessary conditions and exact phase masses}
From the repository root, run
\begin{verbatim}
python evidence/rigidity_20261017/rank6/run_checks.py --repo .
python -O evidence/rigidity_20261017/rank6/run_checks.py --repo .
\end{verbatim}
The two case folders contain their source-bound replay programs and
independent audit receipts. The rank-six $R6\_4\_155$ checker verifies
all $1144$ odd cyclic Clifford decompositions and the all-order path
criterion; the latter does not exclude multiplicity four. The
$R6\_8\_1660$ checker reconstructs every indicator at orders $1$ through
$48$, all $288$ Fourier masses for each conjugate choice, the exact
spin restrictions and the $640$-profile phase Gram domain. Independent
SymPy arithmetic checks the Fourier formulas and the full Ising sequence;
a separate coefficient-convolution implementation exhausts the phase
allocation domain. Explicit rejection controls and both Python execution
modes are included. These checks certify finite arithmetic and necessary
conditions; neither residual has a categorification or a general
nonexistence proof.


\begin{samepage}
None of these necessary alternatives is a categorical model. A claimed
future exclusion must cover the entire necessary domain; a claimed model
must have a categorical construction and complete tensor identification.


\section{Rebuilding the presentation}
After checking the retained evidence, regenerate the ledger exports,
model catalogue, tables and PDFs with
\begin{verbatim}
python software/release.py
python software/model_catalog.py
python software/build_paper_tables.py
python software/build_readme.py
python build_papers.py
\end{verbatim}
Run discovery or regeneration in a working copy, because these commands
write files. The build publishes PDFs only after successful compilation.
All counts must agree with the decision ledger. The supplement retains
accepted proofs, needed dependencies and current residual evidence;
obsolete searches for the three conformally realized rings are omitted.
\end{samepage}

\section{Fermion linking and center indicators}
The manuscript proves the fermion Azumaya linking
identity and checks all $22^4=234256$ instances on the retained
$R6\_4\_155$ based module. This is a necessary compatibility result,
not a construction of module associators. Both rank-six categorification
questions remain unresolved. For $R6\_8\_1660$, an exact character-trace
refinement forces $\nu_3(T)=2$ and leaves two conjugate third-indicator
profiles. A separate exact Verlinde calculation excludes the specified
rank-$37$ witness by forcing a fusion coefficient $1/2$; other profiles
remain open. The source-bound
standard-library replay and proof audit are in
\filename{evidence/unified_20261016/rank6/higher_indicators/}; run its
\filename{verify_r1660_higher_indicators.py} with \verb|--repo .| from the repository
root, also with \filename{-O}. The scoped witness obstruction has its own
standard-library checker
\filename{verlinde_obstruction/verify_r1660_witness_verlinde.py}
under the same evidence directory. It also checks the linear necessary
system, which passes. The fermion identity is reproduced from the
repository root by
\begin{verbatim}
python evidence/unified_20261016/rank6/verify_fermion_crossing.py --repo .
\end{verbatim}
For the complete continuation, run
\begin{quote}\small
\verb|python3 tools/verify_current_publication.py --historical|
\end{quote}
and repeat with \filename{-O}. This mode reconstructs the source-bound
input ledgers required by the certificates before replaying them. It is
a continuation check, separate from the full rank-six campaign described
above.

\end{document}
